\subsection{Formalizing the Connection to Reinforcement Learning}
\label{sec:rl}

Iterative deployment can be interpreted as RL fine-tuning but with the reward
signal left implicit. We prove next that SFT using only valid traces can be seen
as a special case of REINFORCE \cite{williams-jml1992} with an implicit binary reward function.

To show this, we start from the following result:

\begin{proposition}
The update directions of the gradients for SFT using only
valid traces and REINFORCE with binary rewards are identical.
\label{thm:gradient-direction}
\end{proposition}

The proof is included in Appendix~\ref{appendix:proofs}.

In addition to the valid traces generated by the current model, we assume access
to a collection of traces from previous generations of the model. We refer to
the traces produced by current policy $\pi_\theta$ as the on-policy traces, and
those produced by a behavior policy $\pi_\beta$ (earlier generations or any
other external source) as \emph{off-policy} traces.

\begin{proposition}
SFT on a mixture of on- and off-policy valid traces is equivalent to REINFORCE
with binary rewards augmented by importance-weighted contributions from the
behavior policy.
\label{thm:importance-weighted}
\end{proposition}

The proof for Proposition~\ref{thm:importance-weighted} is also included in
Appendix~\ref{appendix:proofs}.

In turn, this proves our original claim that SFT
using only valid traces can be seen as a special case of REINFORCE
\cite{williams-jml1992}.

\begin{proposition}
SFT using only valid traces following the iterative deployment mechanism described in Section~\ref{sec:mechanism} is a special case of REINFORCE with implicitly defined binary rewards.
\label{thm:main-result}
\end{proposition}

\begin{proof}
Follows directly from Proposition~\ref{thm:importance-weighted}.
\end{proof}

